\begin{figure}[!b]
\centering
\includestandalone[mode=buildmissing]{Standalone/Drawings/Firmware_architecture_overview.tex}
    \caption{Firmware module architecture.}
    \label{fig:fw_architecture}
\end{figure}

The firmware is developed in C using the Silicon Labs Simplicity Studio extension for Visual Studio Code, targeting the EFR32MG22E microcontroller with the Simplicity Software Development Kit (SDK) version 2025.12.1. The EFR32xG22 reference manual, together with other Silicon Labs documentation, has been the main reference for this section \cite{silabs_efr32xg22_reference, silabs_efr32xg22_wireless, silabs_Zigbee_EmberZNet_PRO,Simplicity_Studio_6_Overview}. The firmware runs bare-metal, without a real-time operating system.
The SDK uses a cooperative super-loop execution model in which each iteration of the main loop first runs the internal SDK processing via \code{sl\_main\_process\_action()}, followed by the application tick function \code{app\_process\_action()}.
This structure is sufficient because the firmware has a single, well-defined repeating cycle with no concurrent tasks that would require a scheduler.
 
The codebase is decomposed into four application-level modules, each with a clearly bounded responsibility, as illustrated in \cref{fig:fw_architecture}.
The module \code{app.c} owns the top-level state machine and coordinates all other modules.
\code{zigbee.c} handles Zigbee network bring-up, including initial join, stored-network restore, and rejoin fallback.
\code{wheatstone.c} manages the analog front end by controlling bridge excitation through a GPIO output and performing differential IADC sampling.
\code{em4.c} owns the EM4 deep-sleep entry sequence, including BURTC wake-up configuration.
All four modules are built on top of the Simplicity SDK, which provides the Zigbee stack, peripheral drivers (EMLIB), the sleep timer, and the energy management unit (EMU) interface.

The architectural driver behind this design is the battery lifetime requirement.
The system must operate for more than ten years on a primary cell, which demands that both the active time per measurement cycle and the current consumption during deep sleep are minimized.
The firmware therefore structures every cycle around three phases.
The first is a brief active phase in which the Zigbee network is joined or verified, the Wheatstone bridge is sampled, and a Zigbee Cluster Library (ZCL) report is transmitted.
The second is a confirmation phase in which the firmware waits for the transmission callback before committing to sleep.
The third is a deep-sleep phase in EM4 Shutoff, from which the device wakes via a BURTC compare event.
On the EFR32 Series~2, EM4 Shutoff behaves like a full power-on reset, where RAM contents are lost and the chip boots from scratch on wake-up. As a result, each measurement cycle begins identically, which eliminates the need for any persistent state management between cycles and simplifies the firmware considerably.

The resulting firmware image, comprising the four application modules together with the Zigbee stack and the SDK components on which they are built, occupies an address span of \SI{181.7}{\kilo\byte}, equivalent to \SI{35.5}{\percent} of the \SI{512}{\kilo\byte} of flash available on the EFR32MG22E. This substantiates the assessment made during microcontroller selection in \cref{tab:mcu_comparison}, where the comparatively light Zigbee stack was expected to keep the firmware footprint modest, and confirms that the selected device retains substantial storage headroom.
